\chapter{Connections and Transport}
\label{chap:connections-and-transport}

% Closed question: How can information move between histories while preserving
% admissibility?

\begin{chapteressay}

Information moves. A vector at one point on a surface can be
transported to another point along a path; a fact established in
one record can be carried to another record by an admissible
transformation; a theorem proved in one mathematical context can
be applied in another, with appropriate modifications. The
chapter is about how this movement is made admissible.

Transport requires a connection: a structural rule that says how the
moved entity should be modified along the path. Without a
connection, there is no canonical way to compare values at
different points; with a connection, the comparison is well-defined.
The connection is part of the mathematical structure; the transport
is the connection applied.

The chapter develops transport through five mathematical structures.
Parallel transport is the geometric form: a vector on a curved
surface is transported along a curve while maintaining its
orientation relative to the surface. The transport accumulates
holonomy: closed loops produce rotations proportional to the
enclosed curvature. The chapter's example is a vector transported
around a quarter of the globe, accumulating a $90^\circ$ rotation.

The Brouwer fixed-point theorem is the chapter's structural
guarantee. Any continuous self-map of a closed ball has a fixed
point. The theorem is topological: it does not depend on the
specific map, only on the domain's topology. The chapter's
significance is that transport must preserve some structure;
the fixed-point theorem guarantees that at least one structure
survives.

The knowledge partial order is the chapter's example of obstructed
transport. Two incompatible knowledge states have no join in the
lattice; transport across them is obstructed. Quantum mechanics
and general relativity are the canonical instance: each is
admissible alone; their joint extension is not currently
admissible at the Planck scale.

Anderson localization is the chapter's example of disorder-induced
obstruction. An electron in a disordered crystal can be
localized: its wave function does not propagate; the transport
fails. The disorder creates destructive interference that traps
the information. The mathematical analog is the spectral
behavior of random Schr\"odinger operators.

Martin's Axiom is the chapter's extension theorem. For partial
orders satisfying the countable chain condition, locally
consistent choices extend to globally consistent records. The
filter (the global extension) is guaranteed by the axiom. The
chapter's transport is admissible when Martin's condition holds.

The chapter's ground example is the ratio of two counts. When the
counts are at different points on a manifold, comparing them
requires transport. The transport's path-dependence is the
chapter's holonomy. The information is preserved when the path
is admissible; the information is changed when the holonomy is
non-trivial.

The chapter's second concrete example is a poem translated
through three languages and back. The original English is
transported to French, then Japanese, then Russian, then back
to English. The round trip accumulates structural changes:
rhyme is lost, imagery shifts culturally, the tonal register
moves. The final result differs from the original. The
difference is the chapter's accumulated holonomy in literary
form.

The chapter's closed question is:

\begin{center}
\emph{How can information move between histories while preserving
admissibility?}
\end{center}

The chapter's answer: through connections. The connection
specifies the admissible transport. Parallel transport on
manifolds, classical mechanics' Hamilton-Jacobi flow, gauge
fields in physics, and the chapter's typeclass cascades are all
examples of connections. Each has its own admissibility rules;
each accumulates path-dependent residues; each connects local
records into global structures.

Subsequent chapters refine this. Chapter 8 introduces the
calibration that makes multiple observers' records commensurable.
Chapter 9 introduces curvature: the geometric quantity recording
the failure of transport to commute. Chapter 10 introduces
invariance: what survives the transport's relabelings.

\end{chapteressay}

\section{Parallel Transport}
\label{se:parallel-transport}

A vector at the north pole points along a meridian toward Cairo. The
vector is parallel-transported along that meridian to the equator,
maintaining its direction relative to the surface at each step. Upon
reaching the equator, the vector points due south (toward Cairo's
longitude). It is then transported east along the equator, a
quarter of the way around the world, again maintaining its
orientation relative to the surface. The vector now points along the
equator, southward in some sense, but in a different absolute
direction (it has rotated relative to the fixed stars).

Finally, the vector is transported north back to the pole, along the
meridian of its new position. Upon arrival, the vector has rotated
ninety degrees relative to its starting direction. The transport
around the closed triangle (pole, equator-south, equator-east,
back to pole) has accumulated a holonomy of $90^\circ$.

The rotation is not an error. It is the unavoidable consequence of
parallel transport on a curved surface. The earth is curved; the
triangle encloses a spherical area; the holonomy equals the
spherical excess of the triangle (which for a quarter-of-the-globe
triangle is exactly $\pi/2$ steradians, hence $90^\circ$ of
rotation).

The chapter's first claim is that parallel transport carries
information from one point to another along a path, and the result
depends on the path's geometric properties. The vector at the pole
plus the path determines the transported vector at the equator
(and back at the pole). Two different paths can give two different
transported vectors. The difference is the holonomy.

The mathematical formalism is the connection on a vector bundle. A
connection specifies how vectors are transported along curves; it
is a structural object on the manifold (or its tangent bundle).
The Levi-Civita connection on the surface of the earth is the
unique connection compatible with the metric and torsion-free; it
gives parallel transport in the standard geometric sense.

In Lean, the structure appears as the typeclass cascade
\texttt{STEP\_1} through \texttt{ACOLYTE\_PROPAGANDA} in
\texttt{Episode8.lean}. Each step is one transport along the
class stack: a fact established at one level is carried (with
appropriate type adjustments) to the next level. The cascade is
the algorithmic version of parallel transport.

The chapter's claim is that parallel transport is the natural
algebraic operation for carrying information along admissible paths.
The transport is unique once the connection is specified. Different
connections give different transport rules; the choice of
connection is part of the calibration of the admissible structure.

\section{Brouwer Fixed Point}
\label{se:brouwer-fixed-point}

Take a cup of coffee. Stir it. After stirring, some point of the
coffee has ended up exactly where it started. The point may have
traveled in a complicated path during the stirring, but at the end,
some point of the coffee is in its original position.

This is the Brouwer Fixed Point Theorem in its most familiar form.
For any continuous function $f$ from a closed disk to itself, there
exists a point $x$ such that $f(x) = x$. The point is the fixed
point of the function. The stirring is the continuous function;
the fixed point is the unmoved point of the coffee.

The theorem holds in any dimension. For a closed ball in
$\mathbb{R}^n$, any continuous self-map has a fixed point. The
proof requires non-trivial topology (the no-retraction theorem,
which follows from properties of homology or homotopy groups).
The result is widely useful: any equilibrium in a continuous
self-map is a fixed point; the theorem guarantees that at least
one equilibrium exists.

The chapter's second claim is that the existence of fixed points is
forced by topology, not by the specific structure of the map. Any
continuous self-map of a contractible space has a fixed point.
The fixed point is the topological invariant: it exists regardless
of which specific self-map is chosen, as long as the map is
continuous and the domain is contractible.

For transport, the fixed-point theorem ensures that some structure
persists under any admissible transformation. If a record is
mapped to itself by an admissible transport, at least one
component of the record is unchanged. The unchanged component is
the fixed point: the structure that the transport preserves.

\begin{phenom}{Brouwer Fixed Point Theorem}
\label{ph:brouwer-fixed-point}

\PhStatement Let $f : D^n \to D^n$ be a continuous function from the
closed $n$-dimensional ball to itself. Then there exists a point
$x \in D^n$ such that $f(x) = x$.

\PhOrigin Brouwer proved the theorem in 1910 in the context of his
work on the dimension of Euclidean spaces. The theorem is one of
the foundational results of algebraic topology.

\PhObservation For $n = 1$: any continuous function $f : [0, 1] \to
[0, 1]$ has a fixed point. This is a corollary of the
Intermediate Value Theorem applied to $g(x) = f(x) - x$.

\PhConstraint The domain must be a closed ball (or more generally,
homeomorphic to a closed ball: convex, compact, contractible).
The map must be continuous. Without these conditions, the theorem
fails.

\PhConsequence The theorem guarantees the existence of equilibria
in many physical and mathematical systems. It is the foundation
of game theory's Nash equilibrium existence proof, of
differential equations' equilibrium analysis, and of many other
applications.

\end{phenom}

The Brouwer Fixed Point phenomenon is the chapter's instance of a
topological invariant under transport. Any admissible
self-transport of a contractible space has a fixed point; the
fixed point is forced by the space's topology. The transport's
specific rules do not matter; the topology guarantees the
existence.

\section{The Knowledge Partial Order}
\label{se:the-knowledge-partial-order}

Two scientific theories may be incompatible in extreme regimes.
Quantum mechanics and general relativity agree to high precision on
most physical phenomena, but they diverge at the Planck scale:
quantum mechanics predicts fluctuating spacetime at the smallest
scales; general relativity predicts smooth spacetime. The two
theories' predictions at the Planck scale are inconsistent.

The two theories form a partial order on the lattice of possible
physical theories. Each theory is a member of the lattice; the
ordering is by what each theory implies. Quantum mechanics
implies certain predictions; general relativity implies different
predictions. The two theories' \emph{join} would be the theory
implying both wherever they agree: that is, the theory of
non-extreme regimes. The two theories' \emph{meet} would be the
theory implied by both: the theory of the most extreme regimes,
where neither alone is admissible.

The join may not exist in the strict sense. A theory implying both
quantum mechanics and general relativity must reconcile their
predictions at the Planck scale; no such reconciliation is
currently known. The lack of a join is the chapter's
mathematical obstruction to a unified theory.

The chapter's third claim is that knowledge --- in this
philosophical sense --- forms a partial order with potentially
missing joins. Two theories may both be admissible; their joint
implications may not be; the join does not exist in the lattice.
The transport of information across the boundary between the two
theories' regimes is therefore obstructed.

The \texttt{Knowledge.le} relation in \texttt{Episode9.lean}
formalizes this: one knowledge state is at most another when the
first's implications are contained in the second's. The relation
is a partial order; the order has meets and joins (when they
exist); the joins may be missing for incompatible knowledge.

\begin{leanexcerpt}{Knowledge}
inductive Knowledge
  | jarjar : Prop $\to$ Knowledge
  | ledger : Prop $\to$ Fact $\to$ Knowledge $\to$ Knowledge
\end{leanexcerpt}

A \texttt{Knowledge} state is an inductive stack. The
\texttt{jarjar} constructor sits at the base: a single
\texttt{Prop} that the learner has been told (and not yet
checked). The \texttt{ledger} constructor extends a prior
knowledge by attaching a \texttt{Prop}, a \texttt{Fact} that
witnesses it, and the prior stack underneath. The
\texttt{Knowledge.le} relation orders these stacks by
substack-inclusion: one state is at most another when it sits
inside the other's ledger chain.

For two incompatible knowledge states (each consistent
individually, but inconsistent jointly), the join would be the
union of their conclusions; this union violates the consistency
condition; hence the join does not exist in the lattice. The
chapter's claim is that this obstruction is real: not every
pair of records admits an admissible transport.

The lesson: transport is not unrestricted. Information moves
between admissible records only when the records are jointly
admissible. When the joint admissibility fails, the transport is
obstructed; the obstruction is the lattice's missing join.

\section{Anderson Localization}
\label{se:anderson-localization}

In 1958, Philip Anderson studied the behavior of electrons in
disordered crystals. The expected behavior was that electrons would
diffuse through the crystal, conducting electricity in the usual
way. Anderson showed something different. For sufficiently
disordered systems, the electron's wave function does not spread;
it remains localized near its initial position. The disorder
creates destructive interference that traps the wave function.

Mathematically, the Schr\"odinger equation $H \psi = E \psi$ has
solutions that are either extended (spreading throughout the
material) or localized (concentrated in a small region). For
ordered systems, the eigenstates are extended; for disordered
systems above a critical disorder strength, the eigenstates are
localized. The transition is the Anderson transition.

The chapter's fourth claim is that transport can fail. An electron
that is localized cannot propagate; the wave function does not
transport its initial values to other regions of the crystal.
The localization is the obstruction to transport. The chapter's
framework captures this: an admissible record may fail to
transport its information when the underlying medium has too
much disorder.

The mathematical analog of Anderson localization is a random
perturbation of a self-adjoint operator. Let $H_0$ be a
self-adjoint operator with extended eigenstates (e.g., the
Laplacian on a regular lattice). Add a random potential $V$ to
form $H = H_0 + V$. For small $V$, the eigenstates remain
extended; for large $V$, the eigenstates localize.

The phenomenon is universal: any sufficiently disordered system
exhibits localization. The chapter's claim is that this is a
generic feature of transport: too much disorder obstructs the
transport. The obstruction is the chapter's algebraic record of a
non-functioning transport.

In Lean, the analog appears in the typeclass machinery for
\texttt{REPRESENTABLE} in \texttt{Episode3.lean}. A representable
record can be transported through the elaboration cascade; a
non-representable record cannot. The non-representability is
the algebraic version of Anderson localization: the record exists
but cannot be propagated.

The chapter's lesson: not all admissible records are
transportable. Some records are locally admissible but globally
obstructed. The obstruction is a real algebraic phenomenon; the
chapter's discipline carries it honestly as a non-transport.

\section{Martin's Condition and Extension}
\label{se:martins-condition-and-extension}

Martin's Axiom is a set-theoretic principle that strengthens the
Baire Category Theorem. It states: for any partial order $P$
satisfying the countable chain condition, and any countable
family of dense subsets of $P$, there exists a filter on $P$
meeting all the dense subsets. The axiom is independent of
Zermelo-Fraenkel set theory; it is consistent with ZFC and
also independent of the continuum hypothesis.

In the chapter's framework, Martin's Axiom guarantees the
extension of consistent local choices to a globally consistent
choice. The dense subsets are the local conditions; the filter is
the globally admissible record; the existence of the filter is
the chapter's transport guarantee for partial orders satisfying
the countable chain condition.

This is the chapter's fifth claim. Under mild combinatorial
conditions (the countable chain condition), local consistency
admits global extension. A record that is locally consistent at
every point of an admissible family of conditions extends to a
globally consistent record. The extension is the transport across
the family; the existence of the filter is the transport's
guarantee.

The countable chain condition (CCC) is the requirement that any
antichain (a set of pairwise incomparable elements) is countable.
For many natural partial orders (forcing posets in set theory,
the poset of finite functions, etc.), the CCC holds. For some
orders (uncountable antichains exist), the CCC fails; Martin's
Axiom does not apply.

The compiler's analog is the typeclass instance resolution
algorithm. When the elaborator needs an instance, it searches
the global instance database. The search visits dense subsets of
the instance lattice (the candidates of varying specificity);
when a candidate satisfies all the local conditions, the
elaborator selects it. The selection is the chapter's filter:
the globally admissible choice consistent with all the local
requirements.

In Lean, the \texttt{WITNESSED} typeclass in \texttt{Episode10.lean}
formalizes the witness of a transported truth. A
\texttt{WITNESSED} proposition has been transported from some
admissible source through an admissible process; the witness is
the certificate of the transport's success. The certificate is
the filter in Martin's sense: the algebraic confirmation that
the local conditions extend to a global admissible record.

\begin{leanexcerpt}{WITNESSED}
class WITNESSED
    (Value : Type)
    (Carrier : CarrierProcess Value)
    [DISTINGUISHABLE Value Carrier] ... [TRUTH Value Carrier]
  where
  baptism : ReligiousProcess Value Carrier
  witness : Gospel
  risen? : Gospel $\to$ Gospel $\to$ Prop := fun a b =>
    Gospel.le a b $\to$ Gospel.le b a
\end{leanexcerpt}

The chapter closes here. Transport is the algebraic operation of
moving structure between admissible records. The transport is
admissible when the underlying connection is well-defined.
Parallel transport on curved surfaces accumulates holonomy. The
Brouwer fixed-point theorem guarantees the existence of fixed
points under any continuous self-transport. The knowledge
partial order has potentially missing joins; transport across
incompatible knowledge states may be obstructed. Anderson
localization shows that disorder can obstruct transport
algebraically. Martin's Axiom guarantees the extension of
consistent local choices under the CCC.

\begin{bridge}

The chapter's question was how information can move between
histories while preserving admissibility.

The answer is: through a connection that defines admissible
transport. Parallel transport on curved surfaces accumulates
holonomy. The Brouwer fixed-point theorem guarantees the existence
of unmoved points under any continuous self-transport. The
knowledge partial order may have missing joins; transport across
incompatible records is obstructed. Anderson localization shows
that sufficient disorder obstructs transport algebraically.
Martin's Axiom guarantees the extension of consistent local
choices to global admissible records under the CCC.

What remains is the reference structure. Different observers
performing transport may produce different transported values
because they use different connections. The next chapter takes up
calibration: what reference structure lets distinct observers
compare their counts as one physical record?

\end{bridge}

\begin{coda}{Translating a Poem Through Three Languages and Back}

The poem is short. Eight lines. The original is in English, written
by a contemporary American poet now in her seventies. The poem is
well-known: it has been anthologized, taught in workshops, and
translated several times into other languages.

The project this evening is a translation exercise undertaken by a
small literary circle in a coffee shop. The participants are a
poet who reads French, another who reads Japanese, a third who
reads Russian, and a fourth who speaks only English. The
exercise: the French speaker will translate the English original
into French; the Japanese speaker will then translate that French
into Japanese; the Russian speaker will translate that Japanese
into Russian; finally, the English speaker (the fourth participant)
will translate the Russian back into English.

The four-step chain is the chapter's transport in literary form.
The original English is the source; each intermediate version is
a transported translation; the final English is the round-trip
result. The question: what survives?

The first step is the French translation. The French poet reads
the English carefully, several times. She notes the rhyme scheme
(ABABCDCD, with an irregular meter), the central image (a
window at dusk), the recurring word that anchors the poem
(``return''), and the tonal register (intimate, slightly
melancholy). She translates. Her French version preserves the
rhyme scheme (with rhymes appropriate to French phonology), keeps
the central image, finds a French equivalent of ``return'' that
recurs in the same positions, and maintains the tonal register.
The translation takes about forty minutes.

The Japanese translator reads the French. He does not see the
original English. The French version is a poem in its own right;
he treats it as such. He notes the same elements: rhyme,
image, recurrence, tone. Japanese, however, does not have rhyme
in the European sense; Japanese poetic forms use syllable counts
and tonal contrasts instead. The translator adopts the
five-seven-five syllable pattern characteristic of haiku, with
some modification for the eight-line structure. The central image
(the window at dusk) is rendered in Japanese with cultural
attention: a shoji screen rather than a window pane, evening
light rendered with seasonal words. The recurring ``return''
becomes a Japanese word with similar function. The Japanese
translation takes an hour.

The Russian translator reads the Japanese. He does not see the
French or the English. The Japanese poem is now the source. He
notes a different set of elements: the imagery has shifted (shoji
screen, evening seasonal words), the rhythmic pattern is now
syllabic rather than rhymed, the tonal register is more austere
than the original French (and presumably the original English,
which he has not seen). His Russian translation aims for fidelity
to the Japanese; he adopts a Russian rhythmic pattern that
echoes the syllabic count, finds Russian equivalents for the
imagery, and maintains the austere tone. His translation takes
forty-five minutes.

The English speaker now translates the Russian back into English.
She has not seen the original, the French, or the Japanese. She
reads the Russian poem and translates with fidelity to it. Her
English version preserves the Russian imagery, the syllabic-
derived rhythm, the austere tone, the meaning that the Russian
carries.

The final English translation is then compared to the original.
The four participants, the original poet (who joined them by
video), and a few coffee-shop observers look at the two side by
side.

What survives? Some structural features. Both poems are eight
lines. Both have a sense of evening and indoor space. Both
contain a recurring word, though the words are different (the
original's ``return'' has become something like ``come back'' in
the final English, via French \emph{revenir}, Japanese
\emph{kaeru}, Russian \emph{vozvrashchat}'sya). The central
image has shifted: an English window has become a shoji screen
has become a Russian decorative shutter has become an English
shutter. The image is recognizable as ``barrier between inside
and outside, partially transparent''; the specific cultural
form differs.

What has been lost? The rhyme scheme: ABABCDCD is gone, replaced
by free verse. The English language's specific cadence is gone,
replaced by a rhythm influenced by Russian. The intimate
register of the original is muted; the final version is more
austere. The original poet's voice is not recoverable from the
final translation; the final translation has acquired its own
voice.

What has been preserved through transport? The overall
narrative: someone at evening, indoors, contemplating return.
The emotional gesture: a tender remembering. The structural
form: eight lines, evening, recurrence.

This is the chapter's transport in action. The original
information (the poem) has been transported through three
admissible translations and back. Some structural features are
preserved (the eight lines, the emotional gesture, the central
image's category). Other features are lost (the rhyme, the
specific cadence, the intimate register).

The losses are the chapter's holonomy. The transport around the
closed loop (English to French to Japanese to Russian to
English) has accumulated structural changes; the result differs
from the starting point. The accumulation is the chapter's
record of the transport's curvature: the path-dependence of
literary translation across linguistic boundaries.

The discussion at the coffee shop is animated. The original poet
finds the result fascinating. Her recurring ``return'' has
become a different word; her window is now a shutter; her
intimate register has become austere. She does not consider
these losses; she considers them transformations. The poem has
been carried through three languages and has emerged different,
but recognizable.

The exercise ends late. The participants disperse. Each carries
home a small understanding of the chapter's transport
phenomena: information moves between admissible records, but
the path matters, and the round trip never exactly returns.

\end{coda}
